% =====================================================================
\chapter{Vectores, producto escalar y norma}
\label{tema:1}
% =====================================================================

Este tema introduce el espacio $\R^{d}$, sus dos operaciones básicas y las dos herramientas métricas sobre las que descansa todo el módulo: el producto escalar y la norma. La presentación sigue las secciones~1.1 y~1.2 de Strang~\cite{strang} y la clase~1 del curso 18.06SC~\cite{mit1806sc}; las definiciones abstractas de espacio vectorial, norma y producto interior proceden de~\cite[§2.4, §3.1 y §3.2]{mml}, y el punto de vista computacional, de~\cite[cap.~1 y §3.1]{vmls}.

\section{El espacio \texorpdfstring{$\R^{d}$}{ℝᵈ}}

\begin{definicion}[Vector]
Sea $d\ge 1$ un entero. Un \emph{vector} de $\R^{d}$ es una lista ordenada $u=(u_1,\dots,u_d)$ de $d$ números reales, llamados \emph{componentes} de $u$. Dos vectores son iguales si lo son todas sus componentes. El \emph{vector nulo} es $0=(0,\dots,0)$.
\end{definicion}

Un vector de $\R^{2}$ o de $\R^{3}$ admite dos lecturas geométricas equivalentes: un punto del plano o del espacio, o una flecha que parte del origen y termina en ese punto. Para $d>3$ no hay dibujo posible, pero todas las definiciones de este módulo se formulan componente a componente y valen igual en cualquier dimensión.

\begin{definicion}[Suma y producto por un escalar]
Para $u,v\in\R^{d}$ y $c\in\R$ se definen
\[
u+v=(u_1+v_1,\dots,u_d+v_d),\qquad c\,u=(c\,u_1,\dots,c\,u_d).
\]
\end{definicion}

Geométricamente, $u+v$ se obtiene colocando la flecha $v$ a continuación de la flecha $u$ (regla del paralelogramo), y $c\,u$ estira la flecha si $c>1$, la encoge si $0<c<1$ y la invierte si $c<0$. La diferencia $u-v=u+(-1)v$ es la flecha que va del extremo de $v$ al extremo de $u$.

\begin{proposicion}\label{prop:ev}
Con la suma y el producto por un escalar, $\R^{d}$ satisface, para todos $u,v,w\in\R^{d}$ y $a,b\in\R$:
\begin{enumerate}
  \item $u+v=v+u$ y $(u+v)+w=u+(v+w)$;
  \item $u+0=u$ y $u+(-u)=0$;
  \item $a(u+v)=au+av$, $(a+b)u=au+bu$ y $(ab)u=a(bu)$;
  \item $1\,u=u$.
\end{enumerate}
\end{proposicion}

\begin{proof}
Cada igualdad se reduce, componente a componente, a la propiedad correspondiente de los números reales. Por ejemplo, la componente $i$ de $a(u+v)$ es $a(u_i+v_i)=au_i+av_i$, que es la componente $i$ de $au+av$.
\end{proof}

Un conjunto dotado de una suma y de un producto por escalares reales que cumplen las propiedades de la \cref{prop:ev} se llama \emph{espacio vectorial real}~\cite[§2.4]{mml}. Además de $\R^{d}$, son espacios vectoriales el conjunto $\R^{m\times n}$ de matrices (\cref{tema:3}) y el de polinomios de grado a lo sumo $n$. Este módulo trabaja casi siempre en $\R^{d}$.

\begin{definicion}[Combinación lineal]
Sean $v_1,\dots,v_k\in\R^{d}$ y $c_1,\dots,c_k\in\R$. El vector $c_1v_1+\dots+c_kv_k$ es una \emph{combinación lineal} de $v_1,\dots,v_k$ con coeficientes $c_1,\dots,c_k$.
\end{definicion}

\begin{ejemplo}\label{ej:operaciones}
Sean $u=(1,2)$ y $v=(3,-1)$. Entonces $u+v=(4,1)$, $2u=(2,4)$ y $u-v=(-2,3)$. La combinación lineal $2u+v$ vale $(5,3)$.
\end{ejemplo}

Las combinaciones lineales son la operación central del álgebra lineal~\cite[§1.1]{strang}: el producto de una matriz por un vector (\cref{tema:3}) es una combinación lineal de las columnas de la matriz, y la cuestión de qué vectores se obtienen como combinación de otros dados conduce a los conceptos de independencia, base y rango del módulo~3.

\begin{observacion}[Vectores en los modelos de lenguaje]
Un modelo de lenguaje asigna a cada token un vector de $\R^{d}$, su \emph{embedding} (\cref{tema:5}). En el modelo más pequeño de GPT-2 es $d=768$~\cite{radford}, y en GPT-3, con 175\,000 millones de parámetros, $d=12\,288$~\cite[tabla~2.1]{brown}. Anthropic, Google y xAI no publican este valor para los modelos actuales de Claude, Gemini y Grok. Otro ejemplo clásico es el \emph{vector de recuentos} de un documento: fijado un diccionario de $n$ palabras, la componente $i$ es el número de veces que aparece la palabra $i$~\cite[§1.1]{vmls}.
\end{observacion}

\section{El producto escalar}

\begin{definicion}[Producto escalar]\label{def:escalar}
El \emph{producto escalar} (o producto interior euclídeo) de $u,v\in\R^{d}$ es el número real
\[
u\cdot v=\sum_{i=1}^{d}u_iv_i.
\]
Con la notación matricial del \cref{tema:3}, $u\cdot v=u^{\top}v$.
\end{definicion}

\begin{ejemplo}
$(1,2,3)\cdot(4,5,6)=4+10+18=32$. Obsérvese que el resultado es un número, no un vector: el producto componente a componente $(4,10,18)$ es otra operación.
\end{ejemplo}

\begin{proposicion}[Propiedades del producto escalar]\label{prop:escalar}
Para todos $u,v,w\in\R^{d}$ y $a,b\in\R$:
\begin{enumerate}
  \item (simetría) $u\cdot v=v\cdot u$;
  \item (linealidad en cada argumento) $(au+bv)\cdot w=a\,(u\cdot w)+b\,(v\cdot w)$;
  \item (positividad) $u\cdot u\ge 0$, con igualdad si y solo si $u=0$.
\end{enumerate}
\end{proposicion}

\begin{proof}
(i) Se sigue de $u_iv_i=v_iu_i$. (ii) $\sum_i(au_i+bv_i)w_i=a\sum_iu_iw_i+b\sum_iv_iw_i$. Por la simetría, la linealidad vale también en el segundo argumento. (iii) $u\cdot u=\sum_iu_i^{2}$ es una suma de cuadrados, luego es $\ge 0$, y se anula si y solo si cada $u_i=0$.
\end{proof}

Una aplicación $\ip{\cdot}{\cdot}\colon V\times V\to\R$ sobre un espacio vectorial que cumple (i)--(iii) se llama \emph{producto interior}~\cite[§3.2]{mml}. El producto escalar es el producto interior estándar de $\R^{d}$; todos los resultados de este tema, salvo los que usan componentes, valen para un producto interior cualquiera.

\begin{observacion}[Coste]
Calcular $u\cdot v$ en $\R^{d}$ requiere $d$ multiplicaciones y $d-1$ sumas, unas $2d$ operaciones en coma flotante~\cite[§1.5]{vmls}. Un LLM calcula miles de millones de productos escalares por cada token generado, y su organización en productos de matrices (\cref{tema:4}) es lo que permite ejecutarlos eficientemente en una GPU.
\end{observacion}

\section{Normas}

\begin{definicion}[Norma euclídea]\label{def:norma}
La \emph{norma} (o longitud) euclídea de $u\in\R^{d}$ es
\[
\norm{u}=\sqrt{u\cdot u}=\sqrt{u_1^{2}+\dots+u_d^{2}}.
\]
Un vector $u$ es \emph{unitario} si $\norm{u}=1$. Si $u\neq 0$, el vector $\hat u=u/\norm{u}$ es unitario y tiene la misma dirección y sentido que $u$; se dice que $\hat u$ es el \emph{normalizado} de $u$. La \emph{distancia} entre $u$ y $v$ es $\norm{u-v}$.
\end{definicion}

En $\R^{2}$ y $\R^{3}$ la norma es la longitud de la flecha, por el teorema de Pitágoras. Por ejemplo, $\norm{(3,4)}=5$, $\norm{(1,2,2)}=3$ y el normalizado de $(3,4)$ es $(0{,}6;\,0{,}8)$.

\begin{definicion}[Norma]\label{def:norma-general}
Una \emph{norma} sobre un espacio vectorial $V$ es una aplicación $\norm{\cdot}\colon V\to\R$ que cumple, para todos $u,v\in V$ y $c\in\R$~\cite[§3.1]{mml}:
\begin{enumerate}
  \item (definida positiva) $\norm{u}\ge 0$, con igualdad si y solo si $u=0$;
  \item (homogeneidad) $\norm{c\,u}=\abs{c}\,\norm{u}$;
  \item (desigualdad triangular) $\norm{u+v}\le\norm{u}+\norm{v}$.
\end{enumerate}
\end{definicion}

\begin{ejemplo}[Otras normas de $\R^{d}$]
Además de la euclídea, también llamada norma $\ell_2$, se usan la norma $\ell_1$, $\norm{u}_1=\sum_i\abs{u_i}$, y la norma del máximo, $\norm{u}_\infty=\max_i\abs{u_i}$. Para $u=(3,-4)$: $\norm{u}_1=7$, $\norm{u}=5$ y $\norm{u}_\infty=4$. El cuadrado de la norma euclídea de los parámetros es la penalización de la regularización $L_2$ del módulo~8.
\end{ejemplo}

Las propiedades (i) y (ii) de la \cref{def:norma-general} son inmediatas para la norma euclídea. La desigualdad triangular, en cambio, requiere un resultado previo que es probablemente la desigualdad más importante del álgebra lineal.

\begin{teorema}[Desigualdad de Cauchy--Schwarz]\label{teo:cs}
Para todos $u,v\in\R^{d}$,
\[
\abs{u\cdot v}\le\norm{u}\,\norm{v},
\]
y la igualdad se da si y solo si uno de los dos vectores es múltiplo del otro.
\end{teorema}

\begin{proof}
Si $v=0$, ambos miembros valen~0 y $v=0\cdot u$. Sea $v\neq 0$ y sea $t^{*}=(u\cdot v)/\norm{v}^{2}$. Por la positividad y la linealidad del producto escalar (\cref{prop:escalar}),
\[
0\le\norm{u-t^{*}v}^{2}=\norm{u}^{2}-2t^{*}(u\cdot v)+(t^{*})^{2}\norm{v}^{2}
=\norm{u}^{2}-\frac{(u\cdot v)^{2}}{\norm{v}^{2}}.
\]
Multiplicando por $\norm{v}^{2}>0$ se obtiene $(u\cdot v)^{2}\le\norm{u}^{2}\norm{v}^{2}$, y tomando raíces cuadradas, el enunciado. La igualdad se da si y solo si $\norm{u-t^{*}v}=0$, es decir, si $u=t^{*}v$. Recíprocamente, si $u=cv$, entonces $\abs{u\cdot v}=\abs{c}\norm{v}^{2}=\norm{u}\norm{v}$.
\end{proof}

La elección de $t^{*}$ no es arbitraria: $t^{*}v$ es la proyección de $u$ sobre la recta de $v$ (\cref{prop:proyeccion}), y la demostración dice que el cuadrado de la distancia de $u$ a esa recta es no negativo. El ejercicio~\ref{eje:1-cs} propone una demostración alternativa.

\begin{corolario}[Desigualdad triangular]\label{cor:triangular}
Para todos $u,v\in\R^{d}$, $\norm{u+v}\le\norm{u}+\norm{v}$. En consecuencia, la norma euclídea es una norma en el sentido de la \cref{def:norma-general}.
\end{corolario}

\begin{proof}
Por la linealidad del producto escalar y el \cref{teo:cs},
\[
\norm{u+v}^{2}=\norm{u}^{2}+2\,u\cdot v+\norm{v}^{2}\le\norm{u}^{2}+2\norm{u}\norm{v}+\norm{v}^{2}=\bigl(\norm{u}+\norm{v}\bigr)^{2}.
\]
Ambos miembros son no negativos, luego se puede tomar la raíz cuadrada. Las propiedades (i) y (ii) son inmediatas: $\norm{cu}^{2}=\sum_i c^{2}u_i^{2}=c^{2}\norm{u}^{2}$.
\end{proof}

\begin{observacion}
La desigualdad triangular expresa que el camino directo es el más corto: la longitud del lado $u+v$ de un triángulo no supera la suma de las de los otros dos. Aplicada a $u-w=(u-v)+(v-w)$, da $\norm{u-w}\le\norm{u-v}+\norm{v-w}$, que es la desigualdad triangular para distancias.
\end{observacion}

\begin{lecturas}
\item Strang~\cite{strang}, §1.1 (vectores y combinaciones lineales) y §1.2 (longitudes y productos escalares).
\item MIT 18.06SC~\cite{mit1806sc}, unidad~I, sesión «The Geometry of Linear Equations».
\item Deisenroth, Faisal y Ong~\cite{mml}, §2.4 «Vector Spaces», §3.1 «Norms» y §3.2 «Inner Products».
\item Boyd y Vandenberghe~\cite{vmls}, cap.~1 «Vectors» y §3.1 «Norm».
\end{lecturas}

\section*{Ejercicios del tema 1}
\addcontentsline{toc}{section}{Ejercicios del tema 1}

\begin{ejercicio}{1}{eje:1-combinacion}
Sean $u=(1,2)$ y $v=(3,-1)$. Halle los coeficientes $a,b\in\R$ tales que $au+bv=(5,3)$ e interprete geométricamente el resultado.
\end{ejercicio}

\begin{ejercicio}{1}{eje:1-recuentos}
Sea $w\in\R^{n}$ el vector de recuentos de un documento respecto de un diccionario de $n$ palabras que contiene todas las del documento, y sea $\unos=(1,\dots,1)\in\R^{n}$. (a) ¿Qué representa $\unos\cdot w$? (b) ¿Qué significa que $w_{282}=0$? (c) Se define el \emph{vector histograma} $h$, cuya componente $i$ es la fracción de las palabras del documento que son la palabra $i$. Exprese $h$ en función de $w$. (d) Con el diccionario (el, gato, come, pez), en este orden, calcule $w$, $\unos\cdot w$ y $h$ para el documento «el gato come el pez».
\fuente{adaptado de Boyd y Vandenberghe~\cite{vmls}, cap.~1, ejercicio~1.11; se añade el apartado~(d).}
\end{ejercicio}

\begin{ejercicio}{2}{eje:1-normas}
Sea $x=(3,-4,12)$. Calcule $\norm{x}_1$, $\norm{x}$ y $\norm{x}_\infty$, compruebe que $\norm{x}_\infty\le\norm{x}\le\norm{x}_1$ y calcule el normalizado $\hat x$ respecto de la norma euclídea.
\end{ejercicio}

\begin{ejercicio}{3}{eje:1-cs}
Demuestre la desigualdad de Cauchy--Schwarz con el siguiente argumento, distinto del de la demostración del \cref{teo:cs}: para $v\neq 0$, la función $f(t)=\norm{u-tv}^{2}$ es un polinomio de segundo grado en $t$ que no toma valores negativos. Deduzca que $-1\le u\cdot v/(\norm{u}\norm{v})\le 1$ para $u,v\neq 0$.
\end{ejercicio}

\begin{ejercicio}{3}{eje:1-paralelogramo}
(a) Demuestre la \emph{ley del paralelogramo}: $\norm{u+v}^{2}+\norm{u-v}^{2}=2\norm{u}^{2}+2\norm{v}^{2}$ para todos $u,v\in\R^{d}$. (b) Demuestre la \emph{identidad de polarización}: $u\cdot v=\tfrac14\bigl(\norm{u+v}^{2}-\norm{u-v}^{2}\bigr)$. (c) Compruebe con $u=(1,0)$ y $v=(0,1)$ que la norma $\ell_1$ no cumple la ley del paralelogramo, y deduzca que no procede de ningún producto interior.
\end{ejercicio}
